% =====================================================================
%  Capítulo 1 — O problema
% =====================================================================
\chapter{O problema}
\label{cap:problema}
\rubricaJustificativa

\section{Contexto}

Pequenos jardins, hortas em varanda e vasos dentro de casa são cuidados
por pessoas que, em geral, não têm formação agronômica. Elas decidem
quando regar, se devem mover um vaso para perto da janela ou se a planta
está sofrendo com o calor olhando para a planta. Isso tem dois problemas.
Primeiro, quando a planta mostra o sintoma (folha murcha, amarelada ou
queimada), o dano já aconteceu, porque ela responde ao ambiente com horas
ou dias de atraso. Segundo, o que é adequado depende da espécie: a mesma
umidade que faz bem a uma samambaia pode fazer mal a uma suculenta no
vaso ao lado.

Em Teresina a situação é mais difícil. A temperatura fica acima de
\SI{30}{\celsius} em boa parte do ano e, entre agosto e outubro, na
estação seca, a umidade relativa do ar cai bastante. Plantas trazidas de
outras regiões sofrem justamente com essas duas grandezas, e a avaliação
a olho não percebe a mudança a tempo.

A ideia do projeto surgiu de uma conversa de um dos integrantes com a
avó, Teresinha, que cuida de plantas no jardim de casa. Ela relatou duas
dificuldades: lembrar quando foi a última rega de cada planta e saber
quais são as condições adequadas para cada uma. O integrante conhece de
perto a rotina dela e mora perto, o que torna o jardim um lugar acessível
para instalar o nó.

\subsection*{Onde o nó fica, e onde está o primeiro ponto com internet}

Na AP2 o nó saiu da bancada. Ele foi instalado em dois lugares, sempre a
cerca de \SI{1}{\metre} do chão e ligado a um carregador USB de
\SI{5}{\volt} na tomada: num quarto fechado, sobre uma mesa, e num jardim
aberto, fora da casa. A campanha de coleta rodou num quarto, junto a uma janela aberta, com sol direto em parte do dia, com o nó a \SI{50}{\centi\metre} do chão. O primeiro equipamento capaz de receber as
leituras e levá-las adiante é o \emph{central} BLE: na campanha, o
notebook que rodou o gravador, a cerca de \SI{5}{\metre} do nó
(Seção~\ref{sec:caminho}); no produto, o celular de quem cuida do jardim.

A distância máxima com comunicação de sucesso foi de cerca de 20 a
\SI{25}{\metre} dentro da casa, com várias paredes no caminho, e de cerca
de 45 a \SI{50}{\metre} no ambiente aberto, com uma parede só (\CT{21}).
A casa tem WiFi, e o roteador seria o primeiro ponto com internet numa
solução por WiFi, mas o projeto não depende dele: no BLE o roteador nem
entra no caminho do dado.

Daí sai o requisito de alcance (\RNF{06}). Num jardim doméstico, o
celular de quem cuida das plantas fica a no máximo cerca de
\SI{10}{\metre} do nó, quase sempre com uma ou duas paredes no caminho.
O enlace medido alcança o dobro disso dentro de casa, então o produto não
precisa de infraestrutura nova.

\section{Pergunta de projeto}

\begin{quote}
\textbf{Como podemos} informar agricultores domésticos e donos de
jardim, com base em medições de temperatura, umidade e luz,
\textbf{para que} cuidem de cada planta segundo o que a espécie exige,
\textbf{de modo que} saibam quando regar ou reposicionar um vaso antes
de a planta mostrar sinais de dano\textbf{?}
\end{quote}

\section{Quem usa o sistema, e o que muda para essa pessoa}

O usuário é um agricultor doméstico ou dono de jardim, como a Teresinha.
Pode ser alguém que só quer manter as plantas vivas sem muito esforço, ou
alguém que cultiva hortaliças e frutos e quer acompanhar o tempo até a
colheita. Hoje essa pessoa decide pela aparência da planta e pela memória
(\enquote{acho que reguei anteontem}). Com o FitoSfera, ela passa a
contar com medições do ambiente onde a planta está, comparadas pelo
aplicativo com o que a espécie exige, avisos acionáveis
(\enquote{mover a planta Y para um lugar mais iluminado}) e um registro
de cada planta.

Quatro características do produto determinam o que o nó precisa fazer.
É delas que saem os requisitos do Capítulo~\ref{cap:requisitos}:

\begin{table}[H]
\centering
\caption{Do produto ao nó sensor}
\label{tab:produto-no}
\footnotesize
\begin{tabular}{@{}L{4.6cm}L{9.4cm}@{}}
\toprule
\textbf{Característica do produto} & \textbf{Consequência para o nó} \\
\midrule
Monitora umidade, temperatura e luz &
  Três grandezas. Na AP2 a luz deixou de ser claro/escuro e passou a ser
  medida em lux (BH1750), porque o que a espécie exige é uma quantidade
  de luz por dia, e não um estado (\DP{2}). \\
Compara com o que cada espécie exige &
  A comparação mudou para o aplicativo, que conhece a planta
  cadastrada; o nó só mede e entrega (\DP{3}). \\
Fala com o celular de quem cuida do jardim &
  O celular é o central: o nó não precisa de roteador, senha de rede nem
  servidor. O enlace é BLE (\DP{11}). \\
Registra o que aconteceu, com hora &
  O nó não tem relógio: quem carimba a hora é o central, na chegada do
  pacote (\DP{10}). A luz de um período longo sem comunicação não se
  perde, porque o nó a acumula (\DP{8}). \\
\bottomrule
\end{tabular}
\fonte{Elaborada pela equipe.}
\end{table}

As plantas respondem ao ambiente em horas. Na AP1 o nó lia os sensores a
cada \SI{2}{\second}, mais de cem vezes o necessário. Na AP2 o período
passou a \SI{300}{\second} (\DP{9}): continua dez vezes mais rápido que
o fenômeno e reduz a 288 por dia os pacotes que o rádio precisa entregar.

\section{Escopo deste ciclo}

\noindent\textbf{Está neste ciclo (\moduloDoCiclo):}
\begin{itemize}
  \item nó com ESP32 DevKit V1, DHT11 (temperatura e umidade do ar) e
    BH1750 em I\textsuperscript{2}C (luz em lux, até cerca de
    \SI{120000}{lux});
  \item enlace BLE com um serviço próprio: uma característica de leitura
    (notificação de um pacote binário de \SI{19}{\byte}) e uma de
    comando (medir agora e confirmar recebimento);
  \item confirmação de recebimento pela aplicação, luz acumulada no nó e
    identificação de reinícios e de pacotes perdidos;
  \item um gravador BLE em Python, que faz o papel do aplicativo durante
    a campanha e forma o conjunto de dados em CSV;
  \item a campanha de coleta e a análise de qualidade do conjunto de
    dados (\texttt{qualidade.py}).
\end{itemize}

\noindent\textbf{Não está neste ciclo, e por quê:}
\begin{itemize}
  \item \textbf{Aplicativo FitoSfera falando com o nó:} o aplicativo
    ainda não tem o adaptador BLE. Na campanha, o gravador cumpre esse
    papel e segue o mesmo contrato.
  \item \textbf{WiFi, MQTT e \emph{broker}:} a equipe escolheu BLE
    (\DP{11}). Não há \emph{broker}, testamento nem NTP no nó; os
    equivalentes no BLE estão na Seção~\ref{sec:topicos}.
  \item \textbf{Pareamento e segurança:} o nó aceita qualquer central ao
    alcance. Fica para a AP3 (Capítulo~\ref{cap:conclusao}).
  \item \textbf{LEDs de aviso no nó:} saíram junto com o módulo
    fotorresistor. O aviso passa a ser do aplicativo (\DP{3}).
  \item \textbf{Bateria e \emph{deep sleep}:} o nó continua na tomada
    (Apêndice~\ref{ap:energia}).
  \item \textbf{Umidade do solo:} continua sendo a grandeza que melhor
    indica a rega, mas exige uma sonda por vaso.
\end{itemize}

\section{Soluções existentes}

\begin{itemize}
  \item \textbf{Sensores comerciais de planta via BLE} (por exemplo, o
    Xiaomi Flower Care): ficam espetados no vaso e medem umidade do solo,
    temperatura e luz. Usam a mesma tecnologia que o FitoSfera adotou na
    AP2, o que confirma que BLE basta para a distância de uma casa. Mas o
    código é fechado, cada vaso precisa de uma unidade e eles não medem a
    umidade do ar.
  \item \textbf{Aplicativos de lembrete de cuidado} (por exemplo,
    Planta): sugerem rega por calendário e por espécie, mas não medem
    nada. O lembrete é o mesmo na semana mais úmida e na mais seca do
    ano, que é o erro que o FitoSfera quer evitar.
  \item \textbf{Estações de monitoramento e irrigação agrícola:} medem
    com precisão, mas exigem \emph{gateway}, assinatura e instalação que
    não cabem numa varanda.
\end{itemize}

Nenhuma dessas soluções junta, a baixo custo, medições do ambiente e
comparação para cada espécie no celular do usuário. A arquitetura em
camadas usada no Capítulo~\ref{cap:arquitetura} segue
\citeonline{alfuqaha2015}.
